%\textbf{4.2 \emph{Plan Cuadrante} and Spatial Spillovers}

Beyond parallel trends, a causal interpretation of our estimates requires the absence of spillover effects from treated to control municipalities, since such spillovers could bias treatment effects upward. This concern is relevant in our setting. Previous work on deployment-based policing reforms, such as hot-spot policing, shows that crime may be displaced to nearby areas \citep{Blattman_Tobon_JEEA}. Spillovers could also arise through the reallocation of police resources from control to treated municipalities.

In our setting, spillovers are less likely for several reasons. First, as discussed in Section \ref{institutions}, \emph{Plan Cuadrante} introduced beat policing in Chile by reorganizing deployment across the full municipal territory rather than concentrating resources in specific hot spots. Second, Chilean municipalities are large---on average 101,000 inhabitants and 2,000 square kilometers \citep{SINIM}---and criminals typically do not travel such long distances to commit crimes, which makes displacement to other municipalities unlikely \citep{Langella}. Third, and more directly related to police resources, police regulations assign officers to \emph{comisar\'ias}---broader operational units that usually cover an entire municipality---for fixed periods, which limits rapid reallocation of officers across jurisdictions.\footnote{See the Transfer Manual for Carabineros de Chile Personnel available at \url{https://www.carabineros.cl/transparencia/og/pdf/OG_2707_13112019.pdf}}

We nonetheless assess spillovers empirically through two complementary exercises using administrative data at the national level. Online Appendix \ref{sec:appendixC} reports the full results. First, following \citet{crepon2013}, we test whether higher exposure to \emph{Plan Cuadrante} among neighboring municipalities affects outcomes in control municipalities. Online Appendix Figures \ref{fig:figure_continuous_spillovers} and \ref{fig:figure_count_spillovers}---which define neighbor exposure as the share and the number of treated neighbors, respectively---show no meaningful effects across the outcomes we study, suggesting the absence of spatial spillovers. Because exposure to treated neighbors is not randomly assigned, however, this evidence does not fully rule out the possibility that municipalities with many early-adopting neighbors differ systematically from those with fewer such neighbors in other respects that are also correlated with our outcomes.

Second, we estimate the effect of \emph{Plan Cuadrante} on police-reported crime and police resources in untreated municipalities whose neighbors adopt the program. Online Appendix Figure \ref{fig:spillover_admin} summarizes these estimates and again shows no evidence of spatial spillovers.

The exercises above evaluate spatial spillovers using administrative records for the extended municipality sample. We complement this evidence with two additional analyses focused on ENUSC municipalities. First, we re-estimate the administrative-outcome spillover analysis---for reported crime, arrests, police officers, and patrol vehicles---restricting the sample to neighbors of the 46 ENUSC municipalities (Online Appendix Figure \ref{fig:figure_46ENUSC_spillovers_admin}). Second, we test for spillovers directly in ENUSC outcomes. Because ENUSC coverage is limited---46 treated municipalities and 1 never-treated municipality, with only 10 sharing a border and 6 of those treated in different years---we estimate spillover effects from early-adopting ENUSC municipalities onto late-adopting ENUSC neighbors before the latter adopt the program (Online Appendix Figure \ref{fig:figure_results_spillovers_enusc}). Both exercises yield reassuringly null effects.

Taken together, these results indicate that spatial spillovers are unlikely to threaten identification in our setting.

\FloatBarrier